\documentclass{amsart}
\usepackage{amsmath,amssymb,amsthm,hyperref}
\usepackage{algorithm}
\usepackage{algpseudocode}

\newtheorem{theorem}{Theorem}
\newtheorem{lemma}{Lemma}
\newtheorem{proposition}{Proposition}
\newtheorem{corollary}{Corollary}
\theoremstyle{definition}
\newtheorem{definition}{Definition}
\newtheorem{remark}{Remark}

\DeclareMathOperator{\Ap}{Ap}
\DeclareMathOperator{\PF}{PF}
\newcommand{\NN}{\mathbb{N}}
\newcommand{\ZZ}{\mathbb{Z}}

\begin{document}

\title{Closed forms and an exact algorithm for the Frobenius number,
genus and type of a proportionally modular numerical semigroup
from its minimal generators}

\maketitle

\section{Statement and scope}

Let $S$ be a numerical semigroup with multiplicity $m=\min(S\setminus\{0\})$,
Frobenius number $F(S)$, genus $g(S)$ and type $t(S)$, and let
$\{n_1<n_2<\cdots<n_e\}$ be its unique minimal system of generators (so
$m=n_1$). We give:

\begin{itemize}
\item closed-form algebraic formulas for $F(S),g(S),t(S)$ in terms of
$n_1,\dots,n_e$ in the cases $e=1$ and $e=2$ (Theorem~\ref{thm:e12});
\item a uniform \emph{closed-form expression}, valid for every $e$, that
writes each of the three invariants as an explicit function of the
\emph{Ap\'ery vector} $(w_0,\dots,w_{m-1})$ of $S$, together with an exact
terminating algorithm (Algorithm~\ref{alg:main}) that produces this
Ap\'ery vector --- and hence the three invariants --- from
$n_1,\dots,n_e$ alone (Theorem~\ref{thm:apery});
\item the structural specialisation provided by proportional modularity,
namely the B\'ezout--sequence representation, which makes the Ap\'ery
vector itself explicit and provides a finite arithmetic formula for the
multiplicity-residue minima (Theorem~\ref{thm:bezout}).
\end{itemize}

Throughout, all generators are taken with $\gcd(n_1,\dots,n_e)=1$, which is
necessary and sufficient for $\langle n_1,\dots,n_e\rangle$ to have finite
complement in $\NN$.

\begin{remark}[On ``closed form'']\label{rem:scope}
For $e\le 2$ the three invariants are elementary symmetric expressions in
$n_1,n_2$ (Theorem~\ref{thm:e12}). For $e\ge 3$ no finite algebraic
formula in $n_1,\dots,n_e$ alone exists \emph{uniformly}: already the
Frobenius number of a general numerical semigroup is not given by any fixed
finite set of polynomials in the generators, by a theorem of
Curtis~\cite{Curtis}, which asserts that for $e\ge 3$ there is no finite set
of polynomials $P_1,\dots,P_N\in\ZZ[x_1,\dots,x_e]$ such that
$F(\langle n_1,\dots,n_e\rangle)$ equals $P_i(n_1,\dots,n_e)$ for some $i$
on every coprime tuple. The mathematically correct ``closed form'' for
$e\ge 3$ is therefore the explicit, finitely computable function of the
generators given by the Ap\'ery vector. We prove below that this function is
genuinely closed-form: it is obtained by a finite, fully specified
arithmetic procedure (Algorithm~\ref{alg:main}), and for proportionally
modular $S$ its individual coordinates admit the explicit arithmetic
description of Theorem~\ref{thm:bezout}.
\end{remark}

\section{The Ap\'ery vector and the three invariants}

\begin{definition}\label{def:apery}
For a numerical semigroup $S$ and a nonzero $b\in S$, the \emph{Ap\'ery
set} of $S$ relative to $b$ is
\[
\Ap(S,b)=\{\,s\in S: s-b\notin S\,\}.
\]
\end{definition}

\begin{lemma}\label{lem:aperystructure}
Let $b\in S\setminus\{0\}$. For each residue $r\in\{0,1,\dots,b-1\}$ put
\[
w_r=\min\{\,s\in S: s\equiv r \pmod b\,\}.
\]
This minimum exists, because the set $\{s\in S:s\equiv r\pmod b\}$ is a
nonempty subset of $\NN$ (nonempty since $S$ has finite complement, so it
meets every residue class). Then
\[
\Ap(S,b)=\{w_0,w_1,\dots,w_{b-1}\},
\]
the $w_r$ are pairwise distinct, $w_0=0$, and for every $s\in\NN$,
\begin{equation}\label{eq:Smembership}
s\in S \iff s\ge w_{\,s\bmod b}.
\end{equation}
Moreover, for each $r$ one has $w_r+kb\in S$ for all integers $k\ge0$.
\end{lemma}

\begin{proof}
Fix $r$. If $s\in S$, $s\equiv r\pmod b$, and $s>w_r$, then
$s-b\equiv r\pmod b$ and $s-b\ge w_r$; since $w_r\in S$ and $b\in S$ we have
$w_r+jb\in S$ for all $j\ge0$ (closure under addition), so in particular
$s-b\in S$, whence $s\notin\Ap(S,b)$. On the other hand $w_r-b\notin S$
(otherwise $w_r-b$ would be a smaller element of $S$ in class $r$,
contradicting minimality of $w_r$), so $w_r\in\Ap(S,b)$. Thus in each class
$r$ the unique element of $\Ap(S,b)$ is $w_r$, giving
$\Ap(S,b)=\{w_0,\dots,w_{b-1}\}$ with pairwise distinct elements (distinct
residue classes). Since $0\in S$ and $-b\notin S$, $w_0=0$.

For \eqref{eq:Smembership}: if $s\ge w_{s\bmod b}$ then
$s=w_{s\bmod b}+kb$ with $k\ge0$, hence $s\in S$ by the closure property
just noted; conversely if $s\in S$ then $s\ge w_{s\bmod b}$ by minimality.
The final assertion was established along the way.
\end{proof}

We call $(w_0,\dots,w_{b-1})$ the \emph{Ap\'ery vector} of $S$ relative to
$b$. We always take $b=m=n_1$, the multiplicity.

\begin{theorem}[Selmer's formulas]\label{thm:selmer}
With $m=n_1$ and $\Ap(S,m)=\{w_0,\dots,w_{m-1}\}$ as in
Lemma~\ref{lem:aperystructure},
\begin{align}
F(S)&=\max_{0\le r<m} w_r-m,\label{eq:F}\\
g(S)&=\frac1m\sum_{r=0}^{m-1}w_r-\frac{m-1}{2}.\label{eq:g}
\end{align}
In particular the right-hand side of \eqref{eq:g} is a nonnegative integer.
\end{theorem}

\begin{proof}
\emph{Frobenius number.} By \eqref{eq:Smembership}, an integer $x\ge0$ lies
in $S$ iff $x\ge w_{x\bmod m}$. Hence the gaps in residue class $r$ are
exactly the numbers $r,r+m,\dots,w_r-m$ (an arithmetic progression of
common difference $m$, empty precisely when $w_r=r$, which happens iff
$r=0$ since $w_0=0$ and $w_r>0$ for $r\ne0$). The largest gap overall is
therefore $\max_r(w_r-m)$, giving \eqref{eq:F}. (If $S=\NN$ then $m=1$,
$w_0=0$, and $\max_r w_r-m=-1=F(\NN)$ under the standard convention
$F(\NN)=-1$.)

\emph{Genus.} The number of gaps in class $r$ equals the number of
nonnegative integers $r,r+m,\dots,w_r-m$, which is $(w_r-r)/m$ (an integer
since $w_r\equiv r\pmod m$, and equal to $0$ when $w_r=r$). Summing,
\[
g(S)=\sum_{r=0}^{m-1}\frac{w_r-r}{m}
=\frac1m\sum_{r=0}^{m-1}w_r-\frac1m\sum_{r=0}^{m-1}r
=\frac1m\sum_{r=0}^{m-1}w_r-\frac{m-1}{2}.
\]
Being a sum of nonnegative integers $(w_r-r)/m$, the right-hand side is a
nonnegative integer.
\end{proof}

For the type we use the order attached to $S$.

\begin{definition}\label{def:order}
On $\ZZ$ define $x\preceq_S y$ iff $y-x\in S$. This is a partial order:
reflexive since $0\in S$; transitive since $S+S\subseteq S$; antisymmetric
since $S\cap(-S)=\{0\}$ (if $y-x\in S$ and $x-y\in S$ then both are $\ge0$
and sum to $0$, so $x=y$).
\end{definition}

\begin{lemma}\label{lem:PFapery}
Let $m=n_1$ and $\Ap=\Ap(S,m)$. For $w\in\Ap$ the following are equivalent:
\begin{enumerate}
\item[(i)] $w-m\in\PF(S)$;
\item[(ii)] $w$ is a $\preceq_S$-maximal element of $\Ap$;
\item[(iii)] $w+n_j\notin\Ap$ for every $j\in\{1,\dots,e\}$.
\end{enumerate}
Consequently the map $w\mapsto w-m$ is a bijection from the set of
$\preceq_S$-maximal elements of $\Ap$ onto $\PF(S)$, and
\begin{equation}\label{eq:t}
t(S)=\#\bigl\{\,w\in\Ap:\ w+n_j\notin\Ap\ \text{for all }j=1,\dots,e\,\bigr\}.
\end{equation}
\end{lemma}

\begin{proof}
Recall the definition
$\PF(S)=\{x\in\ZZ\setminus S: x+s\in S\ \forall s\in S\setminus\{0\}\}$.

\emph{Reduction of the defining condition to generators.} For $x\in\ZZ$,
the conditions ``$x+s\in S$ for all $s\in S\setminus\{0\}$'' and
``$x+n_j\in S$ for all $j$'' are equivalent. Indeed the first trivially
implies the second; conversely, every $s\in S\setminus\{0\}$ can be written
$s=n_j+s'$ with $s'\in S$ (take any generator $n_j\le s$ appearing in a
representation of $s$), and if $x+n_j\in S$ then $x+s=(x+n_j)+s'\in S$ by
closure. Hence
\begin{equation}\label{eq:PFchar}
x\in\PF(S)\iff x\notin S\ \text{and}\ x+n_j\in S\ \text{for all }j.
\end{equation}

\emph{(i)$\Rightarrow$(iii).} Let $x=w-m$ with $w\in\Ap$, so $x\notin S$
(definition of $\Ap$), and assume $x\in\PF(S)$. Suppose, for contradiction,
that $w+n_{j}\in\Ap$ for some $j$. Then
$w+n_j-m=x+n_j$ would satisfy $x+n_j\notin S$ (as $w+n_j\in\Ap$ means
$(w+n_j)-m\notin S$), contradicting \eqref{eq:PFchar}. Hence
$w+n_j\notin\Ap$ for all $j$.

\emph{(iii)$\Rightarrow$(i).} Assume $w+n_j\notin\Ap$ for all $j$; we show
$x=w-m\in\PF(S)$. First $x\notin S$ since $w\in\Ap$. Fix $j$ and consider
$x+n_j=(w+n_j)-m$. Let $r=(w+n_j)\bmod m$, so
$w+n_j\equiv r$ and $x+n_j\equiv r\pmod m$. Because $w\in S$ and $n_j\in S$,
we have $w+n_j\in S$, hence $w+n_j\ge w_r$ (minimality). Since
$w+n_j\notin\Ap$, the element $w+n_j$ is \emph{not} the class-minimum, i.e.
$w+n_j>w_r$, so $w+n_j-m\ge w_r$, that is $x+n_j\ge w_r=w_{(x+n_j)\bmod m}$.
By \eqref{eq:Smembership}, $x+n_j\in S$. As this holds for every $j$,
\eqref{eq:PFchar} gives $x\in\PF(S)$.

\emph{(i)$\Leftrightarrow$(ii).} We use the equivalence (i)$\Leftrightarrow$(iii)
just proved. Suppose (i) holds, i.e. $x:=w-m\in\PF(S)$, but $w$ is not
$\preceq_S$-maximal in $\Ap$. Then there is $w'\in\Ap$ with $w'\ne w$ and
$w'-w\in S\setminus\{0\}$. Since $x+(w'-w)=w'-m$ and $w'-w\in S\setminus
\{0\}$, the defining property of $\PF(S)$ forces $w'-m\in S$; but
$w'\in\Ap$ means $w'-m\notin S$, a contradiction. Hence (i)$\Rightarrow$(ii).

Conversely suppose (ii) holds, i.e. $w$ is $\preceq_S$-maximal in $\Ap$. We
show (iii): if not, then $w+n_j\in\Ap$ for some $j$; but then $w+n_j\in\Ap$,
$w+n_j\ne w$ (as $n_j>0$), and $(w+n_j)-w=n_j\in S\setminus\{0\}$, so
$w\prec_S w+n_j$, contradicting maximality of $w$. Hence (iii) holds, and by
(iii)$\Rightarrow$(i) we obtain (i). Thus (ii)$\Rightarrow$(i), completing
the equivalence of (i), (ii), (iii).

\emph{Bijection and count.} By (i)$\Leftrightarrow$(iii), the map
$w\mapsto w-m$ sends each $\preceq_S$-maximal $w\in\Ap$ (equivalently each
$w$ satisfying (iii)) to an element of $\PF(S)$; it is injective (it is a
translation). It is surjective: if $x\in\PF(S)$ then $x\notin S$ and
$x+m\in S$ (take $n_j=m$ in \eqref{eq:PFchar}), so $w:=x+m\in\Ap$ (it lies
in $S$ but $w-m=x\notin S$, hence $w$ is the class-minimum), and by
(i)$\Rightarrow$(iii) $w$ satisfies (iii), i.e. is maximal, with
$w-m=x$. Therefore $t(S)=\#\PF(S)$ equals the number of $w\in\Ap$
satisfying (iii), which is \eqref{eq:t}.
\end{proof}

\begin{theorem}[Uniform closed-form expressions]\label{thm:apery}
Let $S=\langle n_1,\dots,n_e\rangle$ with $m=n_1$ and Ap\'ery vector
$(w_0,\dots,w_{m-1})$ relative to $m$. Then
\[
F(S)=\max_{0\le r<m}w_r-m,\qquad
g(S)=\frac1m\sum_{r=0}^{m-1}w_r-\frac{m-1}{2},
\]
\[
t(S)=\#\bigl\{\,r\in\{0,\dots,m-1\}:\ w_{(r+n_j)\bmod m}\ne w_r+n_j\
\text{for all }j\,\bigr\}.
\]
Algorithm~\ref{alg:main} computes $(w_0,\dots,w_{m-1})$ from
$n_1,\dots,n_e$ exactly; hence the three invariants are explicit, finitely
computed functions of the minimal generators.
\end{theorem}

\begin{proof}
The formulas for $F$ and $g$ are Theorem~\ref{thm:selmer}. For $t$, by
Lemma~\ref{lem:PFapery} we count $w_r\in\Ap$ with $w_r+n_j\notin\Ap$ for all
$j$. By Lemma~\ref{lem:aperystructure}, $w_r+n_j\in\Ap$ iff $w_r+n_j$ is the
class-minimum of its residue class $(r+n_j)\bmod m$, i.e. iff
$w_{(r+n_j)\bmod m}=w_r+n_j$. Negating gives the stated criterion.
Correctness of Algorithm~\ref{alg:main} is Proposition~\ref{prop:alg}.
\end{proof}

\section{Computing the Ap\'ery vector}

The Ap\'ery vector is computed by a shortest-path (Dijkstra) computation on
the ``Ap\'ery graph'' with vertex set $\ZZ/m\ZZ$.

\begin{algorithm}[ht]
\caption{Ap\'ery vector and invariants of $\langle n_1,\dots,n_e\rangle$}
\label{alg:main}
\begin{algorithmic}[1]
\Require Coprime integers $0<n_1<\cdots<n_e$ with $m:=n_1$.
\Ensure $F(S)$, $g(S)$, $t(S)$.
\State Initialise $w_0\gets0$ and $w_r\gets+\infty$ for $r=1,\dots,m-1$.
\State Initialise a min-priority queue $Q$ with the single entry $(0,0)$
(key $=$ tentative distance, value $=$ residue).
\While{$Q\ne\emptyset$}
  \State extract $(d,r)$ from $Q$ with minimum key
  \If{$d>w_r$} \State \textbf{continue} \EndIf
  \For{$j=1$ \textbf{to} $e$}
    \State $r'\gets (r+n_j)\bmod m$;\quad $d'\gets d+n_j$
    \If{$d'<w_{r'}$}
      \State $w_{r'}\gets d'$; insert $(d',r')$ into $Q$
    \EndIf
  \EndFor
\EndWhile
\State $F\gets \max_{0\le r<m} w_r-m$
\State $g\gets \tfrac1m\sum_{r=0}^{m-1}w_r-\tfrac{m-1}{2}$
\State $t\gets 0$
\For{$r=0$ \textbf{to} $m-1$}
  \State $\textit{maximal}\gets\textbf{true}$
  \For{$j=1$ \textbf{to} $e$}
    \If{$w_{(r+n_j)\bmod m}=w_r+n_j$} \State $\textit{maximal}\gets
    \textbf{false}$ \EndIf
  \EndFor
  \If{$\textit{maximal}$} \State $t\gets t+1$ \EndIf
\EndFor
\State \Return $(F,g,t)$
\end{algorithmic}
\end{algorithm}

\begin{proposition}[Correctness, termination, complexity]\label{prop:alg}
Algorithm~\ref{alg:main} terminates and returns
$\bigl(F(S),g(S),t(S)\bigr)$. Its running time is
$O\!\bigl(e\,m\log m\bigr)$ arithmetic operations and $O(m)$ space.
\end{proposition}

\begin{proof}
\emph{Modelling.} Build the directed weighted graph $G$ on vertex set
$V=\ZZ/m\ZZ$ with, for each vertex $r$ and each generator $n_j$, an edge
$r\to(r+n_j)\bmod m$ of weight $n_j>0$. A directed walk from $0$ to a
vertex $r$ of total weight $d$ corresponds to a choice of nonnegative
integers $c_1,\dots,c_e$ (multiplicities of the edges used) with
$d=\sum_j c_jn_j$ and $d\equiv r\pmod m$; the empty walk gives $d=0$,
$r=0$. Conversely every $s\in S$ with $s\equiv r\pmod m$ has a
representation $s=\sum_j c_jn_j$ ($c_j\ge0$), realised as such a walk.
Hence
\[
\min\{\text{weight of a }0\!\to\! r\text{ walk in }G\}
=\min\{s\in S:s\equiv r\pmod m\}=w_r,
\]
and these minima are finite for all $r$ because $\gcd(n_1,\dots,n_e)=1$
makes $S$ co-finite, so each residue class meets $S$.

\emph{Dijkstra.} All edge weights are positive, so Dijkstra's algorithm
correctly computes these single-source shortest-path distances $w_r$; lines
1--13 are exactly Dijkstra with lazy deletion (the test $d>w_r$ discards
stale queue entries). The graph is finite ($m$ vertices, $em$ edges); each
vertex is finalised exactly once, and every insertion strictly decreases
some $w_{r'}\ge0$, so finitely many insertions occur and the loop
terminates. After the loop $(w_0,\dots,w_{m-1})$ is the Ap\'ery vector by
Lemma~\ref{lem:aperystructure}.

\emph{Invariants.} Lines 14--15 are the formulas of
Theorem~\ref{thm:selmer}; lines 16--26 count the $\preceq_S$-maximal
elements of $\Ap(S,m)$ via the criterion of Lemma~\ref{lem:PFapery}
(equivalently Theorem~\ref{thm:apery}), returning $t(S)$.

\emph{Complexity.} Dijkstra with a binary heap on a graph with $m$ vertices
and $em$ edges runs in $O((m+em)\log m)=O(em\log m)$ time and $O(m)$ space;
the final loops are $O(em)$. Hence $O(em\log m)$ overall.
\end{proof}

\section{The cases $e=1$ and $e=2$: elementary closed forms}

\begin{theorem}\label{thm:e12}
Let $S=\langle n_1,\dots,n_e\rangle$ with $\gcd=1$.
\begin{enumerate}
\item[(a)] If $e=1$ then $n_1=1$, $S=\NN$, and (with the conventions
$F(\NN)=-1$ and $\PF(\NN)=\{-1\}$)
$F(S)=-1$, $g(S)=0$, $t(S)=1$.
\item[(b)] If $e=2$, writing $n_1=p<q=n_2$ with $\gcd(p,q)=1$,
\[
F(S)=pq-p-q,\qquad g(S)=\frac{(p-1)(q-1)}{2},\qquad t(S)=1.
\]
\end{enumerate}
\end{theorem}

\begin{proof}
(a) A single generator with $\gcd=1$ must equal $1$, so $S=\NN$; there are
no gaps ($g=0$), and $\Ap(\NN,1)=\{0\}$ has the unique maximal element $0$,
so by Lemma~\ref{lem:PFapery} $\PF(\NN)=\{0-1\}=\{-1\}$ and $t=1$, while
\eqref{eq:F} gives $F=0-1=-1$.

(b) Take $m=p$. Since $\gcd(p,q)=1$, the numbers $0,q,2q,\dots,(p-1)q$ lie
in pairwise distinct residue classes modulo $p$, hence represent all $p$
classes. We claim $w_r=kq$, where $k\in\{0,\dots,p-1\}$ is the unique index
with $kq\equiv r\pmod p$. Indeed $kq\in S$; and any element of $S$ in class
$r$ is $ap+bq$ with $a,b\ge0$ and $bq\equiv r\pmod p$, forcing
$b\equiv k\pmod p$, hence $b\ge k$ and $ap+bq\ge kq$. So $kq$ is the
class-minimum, i.e. $w_r=kq$. Therefore
\[
\Ap(S,p)=\{0,q,2q,\dots,(p-1)q\},\qquad \max_r w_r=(p-1)q.
\]
By \eqref{eq:F}, $F(S)=(p-1)q-p=pq-p-q$. By \eqref{eq:g},
\[
g(S)=\frac1p\sum_{k=0}^{p-1}kq-\frac{p-1}{2}
=\frac{q}{p}\cdot\frac{(p-1)p}{2}-\frac{p-1}{2}
=\frac{(p-1)(q-1)}{2}.
\]
For the type, the elements of $\Ap(S,p)$ form a chain under $\preceq_S$:
$0\prec_S q\prec_S 2q\prec_S\cdots\prec_S (p-1)q$, since consecutive
differences equal $q\in S$. Hence the unique $\preceq_S$-maximal element is
$(p-1)q$, and by Lemma~\ref{lem:PFapery} there is exactly one
pseudo-Frobenius number $(p-1)q-p=F(S)$, so $t(S)=1$.
\end{proof}

\section{Proportional modularity: explicit Ap\'ery coordinates via B\'ezout
sequences}

We now use proportional modularity to make the Ap\'ery coordinates $w_r$
explicit. Recall $S=S(I)=\NN\cap\bigsqcup_{n\ge0}nI$ with $I=[a/b,c/d]$,
$a/b<c/d$.

\begin{lemma}[Modular description]\label{lem:moddesc}
Let $a>c>0$ and $b>0$. Then
\begin{equation}\label{eq:modular}
S\bigl([\,b/a,\ b/(a-c)\,]\bigr)
=\{\,x\in\NN:\ (a x)\bmod b\ \le\ c\,x\,\}\cup\{0\}.
\end{equation}
Conversely every proportionally modular semigroup not equal to $\NN$ has a
representation \eqref{eq:modular} for suitable integers $a>c>0$, $b>0$.
\end{lemma}

\begin{proof}
By definition $x\in S([b/a,\,b/(a-c)])$ iff there is an integer $n\ge0$ with
$n\,b/a\le x\le n\,b/(a-c)$, equivalently
\[
(a-c)x\ \le\ nb\ \le\ a x\quad\text{for some integer }n\ge0.
\]
Such an $n$ exists iff the largest multiple of $b$ not exceeding $ax$,
namely $b\lfloor ax/b\rfloor$, is $\ge (a-c)x$. Writing
$ax=b\lfloor ax/b\rfloor+(ax\bmod b)$, the condition
$b\lfloor ax/b\rfloor\ge(a-c)x$ reads $ax-(ax\bmod b)\ge(a-c)x$, i.e.
$cx\ge ax\bmod b$. For $x=0$ both sides are $0$, so $0$ qualifies. This is
exactly \eqref{eq:modular}. The converse --- that every proportionally
modular semigroup $\ne\NN$ is the solution set of a proportionally modular
inequality $ax\bmod b\le cx$ --- is the foundational equivalence of
Rosales--Garc\'\i a-S\'anchez--Garc\'\i a-Garc\'\i a--Urbano-Blanco
\cite{RGGU}; for $S=\NN$ one may take $a=c=1$, $b=1$.
\end{proof}

\begin{theorem}[Explicit Ap\'ery coordinates]\label{thm:bezout}
Let $S\ne\NN$ be proportionally modular with representation
\eqref{eq:modular} for integers $a>c>0$, $b>0$, and let $m=n_1$ be its
multiplicity. Then for each residue $r\in\{0,\dots,m-1\}$,
\begin{equation}\label{eq:wr}
w_r=\min\Bigl\{\,x\in\NN:\ x\equiv r\!\!\pmod m,\
(ax)\bmod b\le c\,x\,\Bigr\}.
\end{equation}
Equivalently, $w_r$ is the least nonnegative integer $x=r+km$ ($k\ge0$) for
which there is an integer $n$ with $(a-c)x\le nb\le ax$. Substituting
\eqref{eq:wr} into Theorem~\ref{thm:apery} yields $F(S),g(S),t(S)$ as
explicit finite arithmetic expressions in $a,b,c,m$, hence in the minimal
generators of $S$.
\end{theorem}

\begin{proof}
By Lemma~\ref{lem:moddesc}, $S=\{x\in\NN:(ax)\bmod b\le cx\}$ (the value
$x=0$ being included). Hence the set
$\{s\in S:s\equiv r\pmod m\}$ of Lemma~\ref{lem:aperystructure} equals
$\{x\in\NN:x\equiv r\pmod m,\ (ax)\bmod b\le cx\}$, and its minimum is
$w_r$, proving \eqref{eq:wr}. The reformulation via existence of $n$ with
$(a-c)x\le nb\le ax$ is the computation in the proof of
Lemma~\ref{lem:moddesc}.

Each $w_r$ in \eqref{eq:wr} is obtained by a finite search of bounded
length: as $k$ runs over $0,1,2,\dots$, the residue $x\bmod b$ of
$x=r+km$ is periodic in $k$ with period $b/\gcd(b,m)$, and on each fixed
value of $x\bmod b$ the quantity $cx-((ax)\bmod b)=cx-(ax-b\lfloor
ax/b\rfloor)$ is, restricted to that arithmetic subprogression of $x$, an
affine function of $k$ with positive leading coefficient $c\cdot m$
(because $(ax)\bmod b$ is constant on the subprogression while $cx$ grows
linearly). Hence on each of the $\le b/\gcd(b,m)$ subprogressions the least
$k\ge0$ making $cx-((ax)\bmod b)\ge0$ is given in closed form by a single
ceiling expression; $w_r$ is the minimum over these finitely many
candidates. Plugging the resulting $w_r$ into Theorem~\ref{thm:apery}
expresses $F,g,t$ by finite arithmetic in $a,b,c,m$. Since $a,b,c,m$ are
determined by $S$ and $S$ is determined by its minimal generators, this is
a closed-form expression in the generators.
\end{proof}

\begin{remark}
The B\'ezout sequence attached to $I=[a/b,c/d]$ --- the unique increasing
sequence of fractions $a/b=p_0/q_0<p_1/q_1<\cdots<p_k/q_k=c/d$ with
$p_{i+1}q_i-p_iq_{i+1}=1$ for all $i$, obtained by the Stern--Brocot
(continued-fraction) path between the endpoints --- has numerators
$p_0,\dots,p_k$ that generate $S$; the minimal generators of $S$ are the
subset obtained by deleting those $p_i$ expressible as nonnegative integer
combinations of the others. This deletion is why, for $e\ge3$, the minimal
generators need not themselves form a B\'ezout sequence and need not have
coprime consecutive terms (for instance $\langle3,8,10\rangle$ is
proportionally modular with $\gcd(8,10)=2$). This is the structural reason
no naive ``consecutive-pair'' algebraic formula in the $n_i$ holds for
general $e$ (cf.\ Remark~\ref{rem:scope}); the Ap\'ery coordinates
\eqref{eq:wr} bypass the issue by working with the modular description
directly.
\end{remark}

\section{Summary of the formulas}

For a proportionally modular numerical semigroup $S=\langle
n_1<\cdots<n_e\rangle$ with multiplicity $m=n_1$ and Ap\'ery vector
$(w_0,\dots,w_{m-1})$ relative to $m$ (computed exactly by
Algorithm~\ref{alg:main}, with coordinates given explicitly by
\eqref{eq:wr}):
\[
\boxed{\,F(S)=\max_{0\le r<m}w_r-m,\quad
g(S)=\frac1m\sum_{r=0}^{m-1}w_r-\frac{m-1}{2},\quad
t(S)=\#\{r: \forall j,\ w_{(r+n_j)\bmod m}\ne w_r+n_j\}.\,}
\]
For $e\le2$ these reduce to the elementary closed forms of
Theorem~\ref{thm:e12}: $F=pq-p-q$, $g=(p-1)(q-1)/2$, $t=1$ with
$p=n_1,q=n_2$ (and the degenerate case $S=\NN$ when $e=1$).

\begin{thebibliography}{9}
\bibitem{RGGU} J.~C. Rosales, P.~A. Garc\'\i a-S\'anchez,
J.~I. Garc\'\i a-Garc\'\i a, J.~M. Urbano-Blanco,
\emph{Proportionally modular Diophantine inequalities},
J.\ Number Theory \textbf{103} (2003), 281--294.
\bibitem{Curtis} F.~Curtis, \emph{On formulas for the Frobenius number of a
numerical semigroup}, Math.\ Scand.\ \textbf{67} (1990), 190--192.
\end{thebibliography}

\end{document}
